% !TeX program = pdflatex
\documentclass[12pt,a4paper]{article}

% Shared preamble for course notes and exercise sets.

\usepackage[T1]{fontenc}
\usepackage{lmodern}
\usepackage[english]{babel}

\usepackage[a4paper,margin=30mm]{geometry}
\usepackage{microtype}
\usepackage{graphicx}
\usepackage{enumitem}

\usepackage{amsmath}
\usepackage{amssymb}
\usepackage{amsthm}
\usepackage{mathtools}
\usepackage{aliascnt}

\usepackage[hidelinks,pdfusetitle]{hyperref}

\numberwithin{equation}{section}
\setlist{itemsep=0.25em,topsep=0.5em}

\theoremstyle{plain}
\newtheorem{theorem}{Theorem}[section]
\newaliascnt{lemma}{theorem}
\newtheorem{lemma}[lemma]{Lemma}
\aliascntresetthe{lemma}
\newaliascnt{proposition}{theorem}
\newtheorem{proposition}[proposition]{Proposition}
\aliascntresetthe{proposition}
\newaliascnt{corollary}{theorem}
\newtheorem{corollary}[corollary]{Corollary}
\aliascntresetthe{corollary}

\theoremstyle{definition}
\newaliascnt{definition}{theorem}
\newtheorem{definition}[definition]{Definition}
\aliascntresetthe{definition}
\newaliascnt{example}{theorem}
\newtheorem{example}[example]{Example}
\aliascntresetthe{example}
\newaliascnt{problem}{theorem}
\newtheorem{problem}[problem]{Problem}
\aliascntresetthe{problem}

\theoremstyle{remark}
\newaliascnt{remark}{theorem}
\newtheorem{remark}[remark]{Remark}
\aliascntresetthe{remark}

\usepackage[nameinlink,noabbrev]{cleveref}

\crefname{theorem}{theorem}{theorems}
\crefname{lemma}{lemma}{lemmas}
\crefname{proposition}{proposition}{propositions}
\crefname{corollary}{corollary}{corollaries}
\crefname{definition}{definition}{definitions}
\crefname{example}{example}{examples}
\crefname{problem}{problem}{problems}
\crefname{remark}{remark}{remarks}

\DeclareMathOperator{\im}{im}
\DeclareMathOperator{\rank}{rank}
\DeclarePairedDelimiter{\abs}{\lvert}{\rvert}
\DeclarePairedDelimiter{\norm}{\lVert}{\rVert}
\DeclarePairedDelimiter{\set}{\lbrace}{\rbrace}


\title{Introduction to University Mathematics}
\author{Joona Piirainen}
\date{Autumn 2026}

\begin{document}

\maketitle

\begin{abstract}
  Notes for the University of Helsinki course MAT11001, taught in Finnish
  from 31 August to 17 October 2026.
\end{abstract}

\tableofcontents

\section{Course information}

The official course information is available on the
\href{https://studies.helsinki.fi/kurssit/toteutus/hy-opt-cur-2627-8d4269d2-7e32-4dce-bcdb-b7fa10fecac6}{University
of Helsinki course page}.

\section{Lecture notes}

\subsection{31 August 2026}

Hello, World!

\begin{definition}
  An integer \(p>1\) is a \emph{prime number} if its only positive divisors
  are \(1\) and \(p\).
\end{definition}

\begin{lemma}\label{lem:prime-divisor}
  Every integer greater than \(1\) has a prime divisor.
\end{lemma}

\begin{proof}
  We use strong induction. Let \(n>1\), and suppose the claim holds for every
  integer strictly between \(1\) and \(n\). If \(n\) is prime, then it divides
  itself. Otherwise, \(n=ab\) for some integers \(1<a<n\) and \(1<b<n\). By
  the induction hypothesis, \(a\) has a prime divisor, which then also divides
  \(n\).
\end{proof}

\begin{theorem}\label{thm:infinitely-many-primes}
  There are infinitely many prime numbers.
\end{theorem}

\begin{proof}
  Suppose instead that \(p_1,\dots,p_k\) are all the prime numbers, and set
  \[
    N=p_1p_2\cdots p_k+1.
  \]
  By \cref{lem:prime-divisor}, some prime \(q\) divides \(N\). Our assumption
  gives \(q=p_i\) for some \(i\), so \(q\) divides both \(p_1\cdots p_k\) and
  \(N\). It must therefore divide their difference \(1\), which is impossible.
\end{proof}

\begin{corollary}
  For every positive integer \(n\), there is a prime number greater than
  \(n\).
\end{corollary}

\begin{proof}
  Otherwise, all prime numbers would belong to the finite set
  \(\{2,\dots,n\}\), contradicting \cref{thm:infinitely-many-primes}.
\end{proof}

\end{document}
